\scanpage{16}
\mspage{8}
Ceci noté, on aura donc, si $j,j'\in J=\operatorname{Ob}\Pi_1$
\[
\begin{aligned}
\operatorname{Isom}_{\Pi_{1,\unsure{\mathrm{é}|\phi}}}(j,j')&\overset{\mathrm{dfn}}=\operatorname{Isom}(\widetilde X'[j],\widetilde X'[j'])\\
&\simeq\operatorname{Isom}_{\pi_1\text{-tors à droite}}(\widetilde X'(j),\widetilde X'(j'))
\end{aligned}
\]
la composition des chemins étant la composition des hom de $\pi_1$-tors à droite.
\deleted{\unsure{Comme}} On a donc reconstruit \inserted{le groupoïde} $\Pi_1$ à l'aide de $\widetilde X'(J)$ et de l'opération libre (à droite) de $\pi_1$ sur $\widetilde X'(J)$.
\deleted{les opérations} Comment reconstituer l'opération de $\mathcal G'$ sur $\Pi_1$ ? L'opération sur $\operatorname{Ob}\Pi_1=J$ n'est autre que l'opération donnée de $\mathcal G'$ sur $J$ (factorisant celle de $\mathcal E'$). Soit $g\in\mathcal G'$, explicitons
\diagram{\begin{tikzcd}[ampersand replacement=\&]
\operatorname{Isom}_{\Pi_1}(j,j') \arrow[r,"g_{\Pi_1}","\sim"'] \arrow[d,equal] \& \operatorname{Isom}(gj,gj') \arrow[d,equal]\\
\operatorname{Isom}_{\pi_1}(\widetilde X'(j),\widetilde X'(j')) \arrow[r,"?","\sim"'] \& \operatorname{Isom}_{\pi_1}(\widetilde X'(gj),\widetilde X'(gj'))
\end{tikzcd}}
Pour ceci, considérons le \inserted{$\pi_1$-}bitorseur
\[
B_g\subset\mathcal E'\qquad\text{image inv. de }g\text{ par }\mathcal E'\mathbin{-}\mathcal G'
\]
On voit que l'on a (grâce à l'op. de $\mathcal E'$ sur $\widetilde X'(I)$]
\[
\begin{aligned}
\widetilde X'(gj)&\simeq\widetilde X'(j)\mathbin{\wedge}_{\pi_1}B_g\\
\widetilde X'(gj')&\simeq\widetilde X'(j')\mathbin{\wedge}_{\pi_1}B_g
\end{aligned}
\]
\mspage{15}
